For an instance of mix with mix formula \(A\) and immediate
subderivations \(D_1\) and \(D_2\), its rank is \(r(A)\) and its height
is \(h(D_1)+h(D_2)\), the sum of the heights of those subderivations.
The reduction is controlled by the lexicographic order on
\[
  \tuple{r(A),h(D_1)+h(D_2)}.
\]
For a cut, the mix instance \(m=n=1\), these measures are its \emph{cut
rank} and \emph{cut height}.

\begin{theorem}[Cut elimination for \(\LK\) and \(\LJ\)]
\label{thm:p07-cut-elimination}
Every \(\LK\)-derivation can be transformed into a cut-free
\(\LK\)-derivation of the same endsequent.  Every \(\LJ\)-derivation can be
transformed into a cut-free \(\LJ\)-derivation of the same endsequent.
\textit{Status: complete proofs in the corresponding annex items.}
\end{theorem}

\paragraph{Complete proof.}
The classical construction is Theorem~\ref{thm:p07-classical-cut-elimination}
and Annex~\ref{annex:p07.classical.cut-elimination}; the intuitionistic
construction is Theorem~\ref{thm:p07-lj-exact-cut-elimination} and
Annex~\ref{annex:p7.lj.exact-cut-elimination}. Each first proves selected-occurrence
Mix admissibility for cut-free premises and then normalizes the original
finite proof tree. The strengthened induction retains unselected copies,
handles contraction and every principal connective/quantifier pair, and
freshens eigenparameters before substitution. Auxiliary alpha-class proofs
are realized with the prescribed literal endsequent. The two constructions
apply respectively to the defined-negation multiset LK and exactly-one-succedent
LJ sheets fixed here; they do not assert the same theorem for equality extensions
or a different primitive-negation presentation.
